\documentclass[10pt]{article} % For LaTeX2e
\usepackage{tmlr}
\usepackage{xcolor}
% For an accepted camera-ready paper, use \usepackage[accepted]{tmlr} instead.
% For a de-anonymized preprint, use \usepackage[preprint]{tmlr} instead.

% V2 tracked changes relative to manuscript_tmlr_revisions.md at
% 384394a113662760cb1485e6ca3e9169c3b79298. Added or revised visible text is blue;
% deletions are omitted, matching the established revision-tracking convention.
\newcommand{\TrackedChange}[1]{\textcolor{blue}{#1}}
\newcommand{\TrackSuppTable}[1]{\textcolor{blue}{#1}}
\newcommand{\TrackSuppTableUnchanged}[1]{#1}
\newcommand{\TrackSuppTableBodyColor}{}
\newcommand{\TrackSuppTableBodyOn}{}
\newcommand{\TrackSuppTableBodyOff}{}

% Optional math commands from https://github.com/goodfeli/dlbook_notation.
\input{math_commands.tex}

\usepackage{amsmath,amssymb}
\usepackage{float}
\usepackage{placeins}
\usepackage{algorithm}
\usepackage{algpseudocode}
\usepackage{graphicx}
\usepackage{booktabs}
\usepackage{array}
\usepackage{longtable}
\usepackage{pdflscape}
\usepackage{hyperref}
\usepackage{url}

\title{FloatSOM: GPU Accelerated, Distributed, Topology-Flexible Self-Organizing Maps}

% Review-stage source: author names, affiliations, emails, and contribution
% statements are deliberately omitted for double-blind review.
\author{}

% TODO(TMLR CAMERA-READY): Replace all three placeholders only for the final
% accepted version; they are unused in anonymous review mode.
\def\month{MM}
\def\year{YYYY}
\def\openreview{\url{https://openreview.net/forum?id=XXXX}}

\begin{document}

\maketitle

\begin{abstract}
GPU-accelerated Self-Organizing Map (SOM) implementations are among the most competitive options for large-scale SOM analysis, but growing dataset sizes increasingly challenge their practical use because workloads no longer fit cleanly within device-memory limits. \TrackedChange{We introduce FloatSOM, a distributed, GPU-accelerated SOM framework. The primary contribution of this work is a scalable systems architecture supporting multi-GPU execution, out-of-memory disk-backed streaming, and hyperparameter tuning. FloatSOM additionally supports scalable training-time graph topologies beyond the regular hexagonal lattice.} \TrackedChange{In our largest benchmark, we trained a 1024-node SOM on a dataset containing 1,000,000,000 samples and 50 features in 6.16 minutes using 8 GPUs across two high-performance-computing nodes.} \TrackedChange{As a secondary evaluation enabled by this framework, we benchmarked graph-based topologies against the conventional hexagonal topology and found that the former yielded a modest reduction in quantization error.}
\end{abstract}

\section{Introduction}

Self-Organizing Map (SOM), originally introduced by Kohonen \citep{kohonenSelforganizingMap1990}, is an unsupervised machine-learning method that uses competitive learning to organize nodes such that they capture the topology of the data. In practice, this topology-preserving representation means SOMs are commonly used to map dataset topology, produce dimensionality-reduced visualizations, and conduct clustering at scale \citep{kangasVariantsSelforganizingMaps1990}. As dataset size and heterogeneity increase, the computational requirements of SOMs grow in both time and memory. Recent tools such as aweSOM provide optimized single-node serial-online SOM execution for clustering analysis \citep{haAweSOMCPUGPUaccelerated2025}. However, many current implementations remain constrained to single-device workloads that must fit within video random-access memory (VRAM, GPU memory), with limited support for distributed compute, out-of-core execution, and modern GPU orchestration.

\TrackedChange{Beyond their scalability constraints, existing implementations are also generally limited in the choice of neighborhood topology they support, restricting SOM training to just conventional rectangular or hexagonal lattices \citep{kohonenSelforganizingMap1990,kohonenEssentialsSelforganizingMap2013}. This is despite a long line of work exploring graph-based SOM topologies as alternatives to the fixed lattice \citep{alahakoonDynamicSelforganizingMaps2000,vasighiDirectedBatchGrowing2017,kangasVariantsSelforganizingMaps1990,spanakisAMSOMAdaptiveMoving2016}. Where graph structures such as minimum spanning trees do appear in deployed SOM workflows, they are typically applied post hoc, derived from a trained lattice map \citep{vangassenFlowSOMUsingSelforganizing2015,jangUseMinimalSpanning2009} rather than as the neighborhood relation that drives learning. Critically, there is currently no implementation of these graph topologies that leverages GPU-accelerated high-throughput processing, leaving them broadly impractical for deployment at scale.}

\TrackedChange{In this paper, we introduce FloatSOM, a distributed, GPU-accelerated SOM framework that removes these constraints. Our primary research question is whether SOM training can be scaled beyond single-device memory and single-GPU throughput through distributed multi-GPU execution and out-of-memory disk-backed streaming. To address this question, FloatSOM implements these systems capabilities together with automated hyperparameter tuning. As a secondary contribution, we also implement scalable minimum-spanning-tree (MST) and relative-neighborhood-graph (RNG) topologies that define the SOM neighborhood during training rather than being derived from a trained map. Accordingly, our primary and secondary contributions are:}

\begin{itemize}
\item \TrackedChange{\textbf{Primary systems contribution:} A distributed, out-of-memory GPU SOM framework that supports multi-GPU execution and disk-backed streaming, scaling SOM training from desktop GPUs to billion-sample workloads across multiple HPC nodes.}
\item \TrackedChange{\textbf{Secondary topology contribution:} Scalable training-time graph topologies (MST and RNG), implemented for the first time in a GPU-accelerated, distributed setting.}
\end{itemize}

\section{Related Work}

Related work on practical SOM deployment spans software implementations, sampling-efficient training, topology design, and model selection.

\subsection{Open-Source SOM Implementations and Systems}

Open-source SOM libraries range from lightweight implementations to systems-oriented packages. MiniSom implements classical serial-online training \citep{vettigliJustGlowingMinisom2018}, while aweSOM provides a Numba-accelerated serial-online SOM for single-node workloads \citep{haAweSOMCPUGPUaccelerated2025}. XPySOM instead provides GPU-accelerated batch training and is the closest executable comparator to FloatSOM's Python/GPU training pathway \citep{manciniXPySomHighPerformanceSelfOrganizing2020}.

Somoclu and GigaSOM provide additional parallel-systems context \citep{wittekSomocluEfficientParallel2017,kratochvilGigaSOMjlHighperformanceClustering2020} and are further elaborated on in Section 8.

\subsection{Sampling Methodologies for SOM Training}

Classical online and batch SOM training schemes traditionally sample every data-point in every training iteration \citep{kohonenSelfOrganizingMaps2001,liuRobustPhenotypingHighly2023}. Consequently, one obvious route to improving SOM speed and scalability is subsampling, which reduces the amount of data used to train the final SOM.

To date, numerous subsampling strategies have been proposed. Beyond naive random sampling, guided subsampling methods include hierarchical dynamic subset selection SOM (HDSSSOM), which concentrates computation on difficult or stale regions of the data \citep{wetmoreSpeedingSelfOrganizingFeature2005}, and adaptive sampling strategies for design-space exploration \citep{itoDesignSpaceExploration2016}. However, systematically benchmarked and openly maintained implementations that compare full-data, random, and guided sampling within a modern high-performance SOM workflow remain limited.

\subsection{SOM Lattice and Adaptive Topologies}

Most practical SOM implementations retain regular rectangular or hexagonal lattices because they simplify neighborhood indexing, visualization, and vectorized updates \citep{kohonenSelforganizingMap1990,kohonenEssentialsSelforganizingMap2013}. Hexagonal lattices are often preferred in the literature and serve as the regular-topology baseline in this manuscript \citep{whiteTopologyMattersNetwork2008,kohonenEssentialsSelforganizingMap2013,forestSurveyImplementationPerformance2020}.

The SOM literature has also explored alternatives to fixed lattices, including dynamic maps and graph-structured neighborhoods \citep{vasighiDirectedBatchGrowing2017,spanakisAMSOMAdaptiveMoving2016,kangasVariantsSelforganizingMaps1990,jangUseMinimalSpanning2009}. MSTs have previously appeared in SOM analyses, but prior uses generally treat the MST as an interpretive structure over an already trained map rather than as the neighborhood relation that drives SOM learning. For example, Jang et al. use MSTs for interpretation, subnode embedding, and map-shape assessment, while FlowSOM overlays an MST on trained SOM codes to visualize relationships among metaclusters \citep{jangUseMinimalSpanning2009,vangassenFlowSOMUsingSelforganizing2015}. This is distinct from the training-time role used here: in FloatSOM, the MST is the operative neighborhood graph during SOM updates, is recalculated from the evolving node-weight geometry, and directly changes the update influence matrix used during learning. Thus, current deployed MST uses in SOM workflows generally configure post hoc connections on trained regular lattice maps, whereas FloatSOM uses MST/RNG graphs as the training neighborhood itself. Earlier SOM variants also discussed MST-defined neighborhoods during learning \citep{kangasVariantsSelforganizingMaps1990}, but these alternatives have not been assessed on large-scale datasets and do not have implementations that are either publicly available or suitable for distributed GPU computation. FloatSOM's topology contribution is therefore a scalable GPU-compatible implementation and large-scale quantification of refreshed graph-based SOM training, rather than a post hoc MST overlay on a conventional trained SOM.

Relative Neighborhood Graphs (RNGs) \citep{toussaintRelativeNeighbourhoodGraph1980} are of particular interest here. To our knowledge, dynamically refreshed RNG neighborhoods have not previously been used as a SOM training topology; we return to the full rationale and implementation for RNG in Section 3.2.2.

\subsection{Hyperparameter Optimization and Fair Comparison}

SOM performance depends strongly on parameters such as map size, initialization, learning-rate schedule, and neighborhood schedule. Prior work shows that these choices can substantially affect observed performance \citep{akindukoSOMStochasticInitialization2016,forestSurveyImplementationPerformance2020}. Because recent SOM system comparisons often emphasize implementation-level performance against existing tools or reference configurations \citep{manciniXPySomHighPerformanceSelfOrganizing2020,kratochvilGigaSOMjlHighperformanceClustering2020}, comparisons based only on untuned settings can be difficult to interpret.

More generally, modern machine-learning workflows increasingly rely on automated hyperparameter optimization rather than manual tuning alone. Tools such as Optuna provide bounded search over large hyperparameter spaces and can support multi-objective optimization, allowing parameter settings to be selected with respect to several benchmark criteria simultaneously rather than collapsed into a single score \citep{akibaOptunaNextgenerationHyperparameter2019}. Accordingly, we test both tuned and untuned reference and FloatSOM implementations to distinguish performance under untuned settings from performance obtained under more suitable hyperparameters.

\section{Methods}

FloatSOM combines four methodological components: sample selection, topology definition, SOM updates, and the compute architecture used to execute them. The framework supports random, full, and HDSSSOM sampling; rectangular, hexagonal, MST, and RNG topologies; and execution modes ranging from local GPU training to single-node and multi-node multi-GPU operation. Hyperparameter optimization is treated as a separate layer applied across these configurations. Figure 1 summarizes this design.

\begin{figure}[ht]
\centering
% SVG source: assets/figures/fig_1.svg
\includegraphics[width=0.5\linewidth]{assets/figures/fig_1.pdf}
\caption{Schematic overview of the FloatSOM methods framing used in this manuscript. Sample selection and topology definition have configurable components that feed into the standard SOM training step, while the compute architecture determines how that same training procedure is executed in practice. The options shown here summarize the FloatSOM configurations discussed in the following subsections.}
\label{fig:1}
\end{figure}

\subsection{Sampling Selector Mathematics}

This section formalizes the sampling policies evaluated in this manuscript.

Let the full dataset be $X=\{x_i\}_{i=1}^{N}$, where $N$ is the total number of samples in the dataset. Let $m$ denote the number of samples presented to the SOM in a given training iteration, or the sampling budget. This budget is either fixed directly or determined as a proportion $\rho$ of the dataset:

\begin{equation*}
m=
\begin{cases}
m_0, & \text{if a fixed budget is specified},\\[4pt]
\max\!\left(1,\left\lfloor N\rho \right\rfloor\right), & \text{if a dataset proportion is used}.
\end{cases}
\tag{1}
\end{equation*}

For index sets $\mathcal{I}_t^{(s)}$ at iteration $t$, the implemented full and random selectors are:

\begin{equation*}
\mathcal{I}_t^{\mathrm{full}}=\{1,\dots,N\},
\qquad
\mathcal{I}_t^{\mathrm{random}} \sim \mathrm{Unif}\!\left(\left\{I\subseteq\{1,\dots,N\}:|I|=m\right\}\right),\ m<N.
\tag{2}
\end{equation*}

If $m\ge N$, random returns the full dataset (no subsampling). Otherwise, to avoid biasing training toward a single fixed subsample, the random selector redraws $\mathcal{I}_t^{\mathrm{random}}$ from the entire dataset at every training iteration. The selected training subset is $X_t^{(s)}=\{x_i:i\in\mathcal{I}_t^{(s)}\}$. Because random subsampling draws uniformly from size-$m$ subsets, we are able to better obtain coverage over the full dataset.

HDSSSOM is implemented in multi-GPU-compatible form as an adaptive sampler that preferentially revisits difficult, under-trained, or stale samples \citep{wetmoreSpeedingSelfOrganizingFeature2005}.

\subsection{Topology Definition}

After initialization, neighborhood relations are defined by the selected topology. For regular-lattice baselines, we support both grid and hexagonal layouts. Consistent with prior SOM guidance, we treat hexagonal as the standard reference topology in this manuscript. We implement the hexagonal lattice topology in accordance with \citep{vettigliJustGlowingMinisom2018}. For MST and RNG, neighborhood structure is derived from the current node-weight geometry using the methodologies below.

Thus, regular-lattice neighborhoods are fixed before training, whereas MST and RNG neighborhoods are recalculated from the evolving node weights during training.

Note that regardless of the selected topology, FloatSOM uses the same node-weight initialization options. Initialization determines only the starting node weights; neighborhood relations are applied afterward according to the selected topology to establish node connections, keeping the initial state comparable across regular-lattice and graph-based runs.

\subsubsection{MST Topology Implementation}

The MST topology replaces fixed lattice neighborhood distance with graph hop distance on a minimum spanning tree built from current SOM nodes. The MST is therefore part of the SOM training rule: it determines which nodes receive neighborhood influence during each topology-refresh interval, rather than serving as a visualization or analysis layer after training. For $P$ nodes in feature dimension $d$, the pairwise node-distance matrix is formed with the standard squared-distance Gram identity, which avoids 3D broadcast tensors and preserves $O(P^2 d)$ dense linear-algebra structure.

In a standard lattice SOM, neighborhood distance is determined by fixed node grid coordinates, which are assigned before training and are independent of the learned node-weight geometry in data space. MST instead uses the node weight vectors to locate the nodes in data space, then calculates the minimum spanning tree between those node locations \citep{kruskalShortestSpanningSubtree1956}. The node distances therefore determine which tree edges are selected, while the neighborhood distance used by the Gaussian update is the graph hop distance along the resulting tree. Topology-derived influence matrices are cached and refreshed according to the topology-refresh schedule.

At the beginning of training, the topology is refreshed every iteration: the previous tree edges are discarded, pairwise distances are recalculated from the updated node weights, and the MST is recomputed from those updated node locations. The refresh rate then decays as training proceeds, so topology updates become less frequent later in training. If the current iteration does not trigger recomputation, the previous graph state and cached influences are reused. Otherwise, the MST is recalculated and the influence cache is rebuilt.

\begin{algorithm}[H]
\caption{Dynamic MST topology update with refresh-triggered recomputation and cached influence reuse.}
\begin{algorithmic}[1]
\State \textbf{Input:} node weights $W_t$, iteration $t$, topology-refresh policy
\State \textbf{Output:} topology state $(E_t, g_t, \text{cached influences})$
\State Query the refresh policy for the current iteration
\If{no topology refresh is due}
\State return previous topology state
\EndIf
\State Compute pairwise squared node distances on GPU
\State Transfer distances to CPU and run Kruskal to obtain MST edges $E_t$
\State Build adjacency from $E_t$
\State Compute all-pairs graph hop distances $g_t$ with chunked GPU Floyd-Warshall
\State Deduplicate active radii and rebuild/update cached influence maps
\State Commit $E_t$, $g_t$, and cache state
\State return topology state
\end{algorithmic}
\end{algorithm}

\subsubsection{RNG Topology Implementation}

Our RNG topology constructs a Relative Neighborhood Graph over current node distances using the standard RNG criterion \citep{toussaintRelativeNeighbourhoodGraph1980}. Unlike MST, RNG is not restricted to a single spanning-tree backbone, so multiple locally supported neighborhood relations can be retained. Conceptually, the absence of this restriction permits the RNG topology to represent both tree-like and denser mesh-like local structures simultaneously throughout the training process. Consequently, we hypothesise that this greater freedom in topology will allow RNG-based SOM maps to better represent a variety of distributions.

\begin{algorithm}[H]
\caption{Dynamic RNG topology update with refresh-triggered recomputation and cached influence reuse.}
\begin{algorithmic}[1]
\State \textbf{Input:} node weights $W_t$, iteration $t$, topology-refresh policy
\State \textbf{Output:} topology state $(E_t, g_t, \text{cached influences})$
\State Query the refresh policy for the current iteration
\If{no topology refresh is due}
\State return previous topology state
\EndIf
\State Compute pairwise squared node distances on GPU
\State Evaluate RNG candidate elimination in chunks using the blocker test
\State Retain surviving RNG edges $E_t$ and build adjacency
\State Compute all-pairs graph hop distances $g_t$ with chunked GPU Floyd-Warshall
\State Deduplicate active radii and rebuild/update cached influence maps
\State Commit $E_t$, $g_t$, and cache state
\State return topology state
\end{algorithmic}
\end{algorithm}

\subsection{Multi-GPU + Larger-Than-Memory Methodology and Implementation}

FloatSOM is implemented for distributed GPU execution, with support for RAM (Random Access Memory, CPU memory) spilling and disk-backed execution when datasets exceed available VRAM and RAM, respectively.

\begin{figure}[tbp]
\centering
% SVG source: assets/figures/fig_2.svg
\includegraphics[width=0.47\linewidth]{assets/figures/fig_2.pdf}
\caption{Multi-GPU data loading and NCCL synchronization schematic. In disk-backed mode, data are sharded to worker-local storage, loaded chunk-wise into pinned host memory, transferred to GPU with transfer overlapped with compute, processed locally, and synchronized by NCCL all-reduce before one weight update per iteration. In RAM mode, data are instead pre-sharded into worker-local CPU RAM and follow the same pinned-memory-to-GPU path without disk reads.}
\label{fig:2}
\end{figure}

\subsubsection{General Multi-GPU Logic}

Distributed execution uses Ray actors with one GPU per worker and NCCL collectives for synchronous aggregation \citep{moritzRayDistributedFramework2018}. At each iteration, every worker processes its assigned shard locally in \texttt{n\_chunks}, accumulates worker-local statistics across those chunks, and synchronizes once per iteration.

For worker $g \in \{1,\dots,G\}$, let $X_{t,g}^{(s)}$ denote the worker-local shard of the selected iteration subset $X_t^{(s)}$. Let $j$ index SOM nodes, let $w_j^{(t)}$ denote the weight vector of node $j$ at iteration $t$, let $b(x)$ denote the best-matching unit (BMU) of sample $x$ under the current weights, let $h_{j,b(x)}^{(t)}$ denote the iteration-$t$ neighborhood influence between node $j$ and the BMU (Best Matching Unit)of $x$, and let $\eta_t$ denote the learning rate at iteration $t$. The local accumulators are:

\begin{equation*}
\begin{aligned}
 U_j^{(g)} &= \sum_{x \in X_{t,g}^{(s)}} \eta_t\,h_{j,b(x)}^{(t)}(x-w_j^{(t)}), \\
 H_j^{(g)} &= \sum_{x \in X_{t,g}^{(s)}} h_{j,b(x)}^{(t)}.
\end{aligned}
\tag{3}
\end{equation*}

$U_j^{(g)}$ is worker $g$'s summed learning-rate-scaled displacement for node $j$, and $H_j^{(g)}$ is the corresponding summed neighborhood influence used to normalize that displacement.

Global synchronized accumulators are:

\begin{equation*}
\begin{aligned}
U_j &= \sum_{g=1}^{G} U_j^{(g)}, \\
H_j &= \sum_{g=1}^{G} H_j^{(g)}.
\end{aligned}
\tag{4}
\end{equation*}

In implementation, $U_j^{(g)}$ and $H_j^{(g)}$ are accumulated across the worker's \texttt{n\_chunks} and synchronized once per iteration. After synchronization, $U_j$ and $H_j$ define the global update numerator and normalization denominator for node $j$. Under weighted normalization, the normalized update is given by $U_j/H_j$ up to numerical safeguards; optional momentum is then added to this normalized update before applying the resulting change to the node weight $w_j^{(t)}$.

For $P$ nodes and $d$ features in float32, the two update accumulators contain $P d+P$ values, so their logical all-reduce payload is $4P(d+1)$ bytes per worker per iteration, plus one 4-byte sample count. A ring all-reduce sends and receives approximately $2(G-1)/G$ times that payload per worker for $G$ workers. This communication is independent of the number of observations $N$; increasing $N$ instead increases worker-local computation and data staging. With observation chunks of at most $C$ rows, the dominant bounded worker arrays scale as $O(Cd+Pd)$ rather than requiring the full $O(Nd)$ dataset in GPU memory. Graph-topology distance and influence structures add both computation and storage terms that scale with $P^2$. Their construction and refresh therefore become more expensive as node count increases, particularly for RNG. FloatSOM tiles these structures and can spill them from VRAM to system RAM when necessary. Spilling is a memory-management response that may add transfer overhead, but it is not the sole source of the topology-dependent runtime increase.

\subsubsection{Multi-GPU Implementation Details}

As shown in Fig. 2, each worker uses a chunked loading path from CPU memory to GPU memory. In streaming mode, data are distributed to worker-local disk shards and then read chunk-by-chunk into pinned host memory before transfer to GPU; in RAM mode, data are pre-sharded into worker-local CPU RAM and follow the same pinned-memory path without disk reads. FloatSOM overlaps data transfer and compute across CUDA streams, uses JIT-compiled kernels (Just-In-Time) for the core BMU and update steps, and synchronizes worker-local accumulators once per iteration with NCCL all-reduce. Weights therefore remain resident on worker GPUs across iterations. For large topologies that would otherwise exceed VRAM, FloatSOM also supports spilling the associated graph-distance or influence structures from VRAM to system RAM.

\subsubsection{Larger-Than-Memory Capable Topology Updates}

Additional larger-than-memory support for topology updates is provided through topological chunking. Topological chunking applies the same idea to topology-side computations within each worker. When graph-distance or influence structures would otherwise exceed a worker's VRAM, those computations are tiled and evaluated in bounded pieces rather than materialized at once. This topology-side chunking is not depicted in the Fig. 2 data path, but it follows the same per-worker bounded-memory execution rule.

\subsection{Multi-Objective Hyperparameter Optimization}

We use Optuna as an automated multi-objective hyperparameter optimization tool to derive near-optimal performance and corresponding hyperparameters for each FloatSOM configuration under a given sampling and topology combination \citep{akibaOptunaNextgenerationHyperparameter2019}.

\section{Experimental Setup}

We use two benchmark protocols: an Optuna quality benchmark and a speed scaling benchmark. The first evaluates algorithmic quality and tuned attainable performance, and the second evaluates runtime and distributed scaling behavior. All production Optuna and speed benchmarks reported in this manuscript were executed on Gadi at the National Computational Infrastructure (NCI), Australia, on gpuvolta nodes \citep{HPCSystemsNCI}, using a consistent multi-GPU environment across runs.

\subsection{Optuna benchmark protocol}

We use Optuna-based multi-objective optimization to determine the best attainable performance and corresponding hyperparameters for each sampling (full and random) and topology (hexagonal, MST, and RNG) combination, optimizing jointly for train and holdout quantization error ($QE_T$ and $QE_H$). In this protocol, each dataset-topology-sampling configuration is optimized for 200 trials across 10 seeds under variant-appropriate search constraints. The reported runs used the standard in-memory batch path rather than the Ray-distributed execution stack. This corresponds to:

\begin{equation*}
14_{\text{Datasets}} \times 10_{\text{Seeds}} \times 3_{\text{Topologies}} \times 2_{\text{SamplingMethods}} \times 200_{\text{Trials}}= 168{,}000 \text{ Runs}
\end{equation*}

Map size was fixed rather than optimized: every Optuna quality run used $P=100$ nodes, implemented as a $10\times10$ lattice for the hexagonal topology and as 100-node graphs for MST and RNG.

The neighborhood-radius search space was shared across topology families. In particular, \texttt{initial\_radius} was an Optuna-optimized parameter for hexagonal, MST, and RNG runs, with the same search interval of 0.5 to 10.0 in each case. Thus, the hexagonal topology comparisons below use a tuned hexagonal comparator rather than a default-radius hexagonal baseline.

Sampling comparisons in Section 5.2 compare full versus random paired analyses pooled across all topologies. Topology comparisons in Sections 5.3-5.4 use full sampling runs across hexagonal, MST, and RNG. Otherwise, the remaining analyses use full sampling. A separate focused HDSSSOM pilot is reported in Section 5.2 and summarized in Supplementary Table S2. This pilot additionally included the HDSSSOM sampling methodology alongside the full and random sampling methodologies.

\subsubsection{Optuna benchmark datasets and preprocessing}

The Optuna benchmark uses 14 synthetic and real scikit-learn datasets spanning low- and high-dimensional regimes \citep{pedregosaScikitlearnMachineLearning2011}; dataset metadata are listed in Supplementary Table S1. Synthetic datasets are generated according to the random seed, whereas real datasets are loaded directly from sklearn. Across the Optuna protocol, inputs are standardized feature-wise to zero mean and unit variance using \texttt{StandardScaler}, then deterministically permuted and split into fixed 70/30 train-holdout partitions. We retain both partitions to evaluate both training-set representation and transfer to unseen samples.

\subsubsection{Optuna quality metrics}

Our primary quality metric is Quantization Error ($QE$) \citep{kohonenSelfOrganizingMaps2001}. We report $QE$ on both the training partition ($QE_T$) and the holdout partition ($QE_H$). The training value reflects how well the fitted SOM represents the samples it was trained on, whereas the holdout value evaluates the same trained map after projecting previously unseen samples onto it. The holdout partition is therefore not used during fitting; it is used only to evaluate the trained map, analogous to a test set in conventional model-training workflows.

Balanced $QE$, denoted $QE_B$, is defined as the mean of $QE_T$ and $QE_H$. $QE_B$ is therefore a composite endpoint that weights data representation fidelity (train) and generalizability (holdout) equally. Note that in the Optuna runs, $QE_T$ and $QE_H$ are optimized jointly as a two-objective vector, with $QE_B$ only calculated \emph{post hoc}.

We quantify local BMU-neighborhood ordering using Mean Tied Rank (MTR), following its application to comparisons among alternative SOM grid topologies \citep{lopez-rubioGridTopologiesSelforganizing2014}. For each sample $x_i$, let $b_i^{(1)}$ and $b_i^{(2)}$ denote the first and second best-matching units (BMUs), respectively. Starting from $b_i^{(1)}$, we group all non-winning units by the number of hops along the shortest path in the SOM graph. Let $d_i$ be the number of hops from $b_i^{(1)}$ to $b_i^{(2)}$, let $S_{i,d_i}$ be the set of units exactly $d_i$ hops away, and let $L_{i,d_i}$ be the number of units fewer than $d_i$ hops away. Because every unit in $S_{i,d_i}$ is tied at the same hop distance, the second BMU is assigned their average ordinal rank:

\begin{equation*}
\begin{aligned}
\tau_i
&= \underbrace{L_{i,d_i}}_{\text{units fewer hops away}}
 + \underbrace{\frac{|S_{i,d_i}|+1}{2}}_{\text{average rank at the same hop distance}}, \\
\mathrm{MTR}
&= \underbrace{\frac{1}{N}\sum_{i=1}^{N}\tau_i}_{\text{mean tied rank across the }N\text{ evaluated samples}}.
\end{aligned}
\end{equation*}

Hexagonal, MST, and RNG maps define adjacency differently: hexagonal connections are fixed by the lattice, MST connections form a sparse tree derived from the node weights, and RNG connections are also weight-derived but can give nodes different numbers of neighbors. MTR applies the same graph-distance ranking procedure to each trained map while accounting for topology-specific differences in node degree and tied hop distances. The $L_{i,d_i}$ term counts all units fewer hops away, while the second term assigns the average rank across the $|S_{i,d_i}|$ units at the same hop distance as the second BMU. Lower MTR indicates that first and second BMUs tend to occur earlier in this graph-specific ordering, but the raw scale remains graph-specific because shell sizes and adjacency construction differ among topology families.

For MST and RNG, adjacency is derived from the same prototype-vector distances that determine first and second BMU identity, so we additionally evaluated MTR using a per-map fixed-adjacency permutation null. For each fitted map and evaluation split, prototype identities were randomly permuted among the $P=100$ nodes 1,000 times while retaining the final graph connections, and MTR was recomputed. We report the observed-to-null ratio $R_{\mathrm{MTR}}=\mathrm{MTR}_{\mathrm{observed}}/\overline{\mathrm{MTR}}_{\mathrm{null}}$; the complete procedure is provided in Supplementary Methods S2.

We additionally report node utilization, defined as the fraction of nodes selected as a BMU, and its complement, dead-node fraction. These post hoc diagnostics are computed on training and holdout splits and balanced as for $QE$.

\subsubsection{XPySOM calibration protocol}

We calibrated the FloatSOM batch pathway against XPySOM using the 14 benchmark datasets, 10 shared seeds, full sampling, a $32\times32$ map, and one run per dataset--seed condition without top-$k$ selection. Both implementations used the same deterministic 70/30 train--holdout split, number of batch epochs, and XPySOM-like untuned hyperparameters. XPySOM's initialized node weights were transferred into FloatSOM after reconciling the two implementations' node-coordinate ordering, so paired runs began from corresponding prototypes. The primary implementation-concordance comparison used the hexagonal topology in both implementations; the associated MST and RNG FloatSOM paths retained the same XPySOM hexagonal reference to quantify the additional effect of changing the FloatSOM training topology.

\TrackedChange{We recorded train, holdout, and balanced $QE$ together with training time. Within each dataset, paired percentage effects across the 10 seeds were summarized using the Wilcoxon signed-rank location estimate, its exact distribution-free 95\% interval, and a two-sided paired Wilcoxon test; global rows pool the 140 dataset--seed pairs. Supplementary Fig. S1 and Supplementary Table S4 report the hexagonal implementation-calibration results; the detailed MST and RNG comparisons are available in the accompanying Zenodo archive.}

\subsubsection{aweSOM quality baseline}

We evaluated aweSOM as an external quality baseline on Iris, Wine, Digits, Breast Cancer, and Olivetti Faces using 10 shared seeds. Preprocessing and the deterministic 70/30 train--holdout split followed Section 4.1.1. Both implementations used 100-node maps and 50 matched data passes. FloatSOM used a $10\times10$ hexagonal lattice and the untuned XPySOM defaults employed elsewhere in the manuscript, while aweSOM used its native $10\times10$ rectangular lattice with Chebyshev neighborhood distance and its native default learning and initialization settings. We calculated train, holdout, and balanced $QE$ and performed paired comparisons within dataset and seed. Dataset-level p-values were Holm-adjusted separately within each $QE$ endpoint.

\subsection{Speed scaling benchmark protocol}

The speed scaling benchmark evaluates runtime and distributed scaling behavior in different compute and algorithm configurations. Speed scaling is evaluated with runs across $G\in\{1,2,4,8\}$ GPUs under a series of fixed scaling protocols. Runtime summaries are computed from repeated executions per configuration and reported as both absolute training time and efficiency relative to the corresponding 1-GPU comparison. We compute scaling efficiency as $E_G=(T_1/T_G)/G \times 100\%$, where $T_1$ is the 1-GPU runtime and $T_G$ is the runtime on $G$ GPUs. A value of 100\% corresponds to ideal linear scaling; lower values indicate that communication, orchestration, or I/O overhead has reduced the realized speedup, while values above 100\% can occur when additional GPUs also change the memory or data-staging regime. These scaling runs use the Ray-orchestrated distributed execution layer built on top of the standard FloatSOM training path \citep{moritzRayDistributedFramework2018}. Accordingly, the scaling figures in Sections 6.1-6.2 and the runtime/scaling comparison reported later against XPySOM should be interpreted as distributed-execution results rather than the in-memory Optuna path.

Some 1-GPU baselines were unavailable because the corresponding runs did not complete within the fixed benchmark timeout, particularly for workloads entering disk-backed execution. Because efficiency is defined relative to the 1-GPU runtime, we estimated these missing denominators by local extrapolation from the last successful 1-GPU point on the same curve. These extrapolated denominators are plotting references rather than measured runtimes and are interpreted only descriptively.

Topology speed comparisons include hexagonal, MST, and RNG, with harmonized workload settings so ratios isolate topology-associated runtime effects. A dedicated random versus full comparison, run on $G\in\{1,2,4\}$ GPUs, is additionally included to isolate sampling-specific runtime effects independently of the multi-GPU scaling runs.

XPySOM was selected as the executable external baseline because it most closely matches FloatSOM's Python/GPU batch-training regime. We also attempted aweSOM on the standard speed workload, but all five runs reached the 1800-s timeout before completing 60\% of $N$ online updates. Matching FloatSOM's 10 batch iterations would require $10N$ pointwise updates, so aweSOM was not included in the timed comparison.

\subsubsection{Scaling benchmark datasets}

The speed benchmark uses synthetic random matrices with uniform values in $[0,1]$. No train-holdout split is used here because the objective is runtime rather than generalization, so all generated data are used for training. We use a shared default configuration of $10^7$ samples, 50 dimensions, a $32 \times 32$ SOM grid, and 10 training iterations, and then vary one workload axis at a time: sample count ($10^6$--$10^9$), feature dimension (50--5000), or grid side length (8--64). The exact axis values used for each scaling panel are provided in Supplementary Table S3.

Additional CPU and RAM resources attached to each GPU are scaled linearly with GPU count while keeping the software environment and benchmark procedure consistent across runs. For scaling benchmarks, each run is timed out if it exceeds 30 minutes of wall time. Per-GPU resource allocation is fixed at one NVIDIA V100 (32 GB VRAM), 12 CPU cores, 90 GB system RAM, with one full node comprising 4 GPUs and their associated CPUs, and with 400 GB of associated local disk storage.

\subsection{Hyperparameter Tuning and Stability}

To quantify parameter-tuning benefit, we performed an explicit paired analysis between the tuned configuration and the XPySOM untuned default reference. For each sampling mode and topology combination, we first extracted the parameter settings from the best-performing Optuna runs under the benchmark objective for that combination. We then distilled these per-seed best-performing tuned settings into deployable default configurations by taking the mean of numeric parameters and the mode of categorical parameters.

The stability analysis focuses on four tuned hyperparameters that govern SOM training dynamics. These settings determine how the map is initialized and how updates evolve over training: the initial radius sets the early neighborhood scale around each BMU, the initialization method sets the starting node weights, the radius decay type controls the shift from broad global organization toward local refinement, and the momentum-use parameter determines whether each update retains part of the previous update direction.

\subsubsection{Tuned Configuration versus Untuned Reference Analysis}

Fixed configurations were rerun under matched dataset, seed, topology, and split keys. All tuned and untuned diagnostic reruns retained this same fixed $P=100$ node count across the three topologies. The final topology diagnostic comprised all 14 benchmark datasets, 20 shared random seeds, full sampling, and each of the hexagonal, MST, and RNG topologies, giving 280 dataset--seed units per paired contrast. Preprocessing and the deterministic 70/30 train--holdout split followed Section 4.1.1. The untuned profile used the same XPySOM-like reference settings for all topologies, whereas the tuned profile used the fixed topology-specific full-sampling configurations derived above. Deployment analyses compare these tuned FloatSOM configurations with untuned hexagonal XPySOM, while topology diagnostics compare tuned and untuned FloatSOM within and between topologies. These reruns estimate deployable fixed-configuration performance; the top-$k$ Optuna summaries instead estimate attainable performance within the search budget.

\TrackedChange{For Supplementary Fig. S3, tuned and untuned values were paired within dataset, seed, topology, and split. Supplementary Table S7 pools the 280 matched dataset--seed differences for each prespecified profile--topology--metric contrast and reports the mean paired effect, two-sided one-sample paired $t$-test confidence interval, Cohen's $d_z$, raw p-value, and win/loss/tie counts. These raw p-values are descriptive diagnostics rather than a multiplicity-adjusted family; effect sizes and confidence intervals are the primary summaries. The corresponding dataset-level mean observed metrics are available in the accompanying Zenodo archive. Dataset-level forest-plot results use the convention described above, and overall rows pool all matched pairs with the same paired $t$-test and confidence-interval calculation.}

\subsubsection{Hyperparameter stability and dataset-type stratification}

Hyperparameter stability is important because it determines whether a method can be deployed reliably without bespoke retuning. We therefore extracted the top-ranked tuned Optuna trial separately for each seed and compared the recovered hyperparameter values across seeds. We define hyperparameter stability as the variation in these recovered settings. For numeric parameters, stability is measured by the relative difference between runs, $|a-b| / \max(|a|,|b|,\varepsilon)$, where $a$ and $b$ are the values of a given parameter for the two compared seeds and $\varepsilon=10^{-12}$. These relative differences are then averaged across numeric parameters, with lower values indicating higher stability. For categorical parameters, stability is defined as the mismatch rate across the same seed pairs. We report both mean per-parameter stability and an equal-weight overall stability summary across datasets.

\subsubsection{Matched initial-radius sensitivity analysis}

To determine whether topology differences could be explained by topology-specific treatment of neighborhood radius, we performed a controlled initial-radius sweep. Hexagonal, MST, and RNG maps were evaluated at $r\in\{0.5,0.75,1.0266,1.5,2,3,5\}$ using the same 20 random seeds on each of the 14 benchmark datasets. The value $r=1.0266$ was the independently Optuna-selected hexagonal full-sampling radius and is denoted \texttt{Hex optimal} in Fig. 9; this label does not denote the minimum observed for hexagonal within the fixed-configuration sweep. Apart from topology and the deliberately varied initial radius, the training configuration was fixed: full sampling, random initialization, asymptotic radius decay, momentum enabled with initial momentum 0.6069, and XPySOM-compatible normalization. Runs were therefore matched by dataset and seed, ensuring that topology families received the same radius values rather than different search ranges or topology-specific radius schedules.

\TrackedChange{We evaluated balanced $QE$, balanced MTR, and balanced node utilization. Because absolute $QE$ scales differ substantially among datasets, each dataset--seed $QE_B$ value was divided by the matched hexagonal $QE_B$ at $r=1.0266$. Ratios were averaged on the log scale across seeds within each dataset and then across datasets with equal dataset weight; exponentiating this mean yields the geometric-mean ratio shown in Fig. 9A. Supplementary Table S9 reports the exact normalized $QE_B$ values and confidence intervals for every topology and radius. Balanced MTR and balanced node utilization are retained in their observed units in Fig. 9B and Fig. 9C, respectively. Along each topology's response curve, direction-aware percentage changes from $r=1.0266$ were calculated for all 280 matched dataset--seed units and tested against zero using two-sided one-sample $t$-tests; Benjamini--Hochberg adjustment was applied across the 54 non-anchor radius--topology--metric tests. These tests assess within-topology radius response.}

\TrackedChange{We additionally conducted post hoc cross-topology $QE_B$ contrasts at each radius. For every topology pair, dataset, seed, and radius, we calculated the log ratio of the comparator $QE_B$ to the reference $QE_B$, averaged log ratios across the 20 seeds within each dataset, and applied a two-sided one-sample $t$-test to the 14 equally weighted dataset effects. Exponentiated means and confidence limits are geometric-mean $QE_B$ ratios, with values below 1 favouring the comparator. Benjamini--Hochberg correction was applied across all 21 topology-pair--radius contrasts. Supplementary Table S10 reports the estimates, intervals, raw p-values, and adjusted q-values.}

\TrackedChange{We then identified the per-topology minimum $QE_B$ across the seven tested radii using the same dataset-balanced geometric-mean summary as Fig. 9A. Paired tests compared $QE_B$ between the configurations at those topology-specific radii, with Benjamini--Hochberg correction applied across the three topology contrasts (Supplementary Table S11).}

\subsubsection{Execution-path concordance analysis}

\TrackedChange{To determine whether the quality diagnostics depended on execution pathway, we compared local CuPy and Ray streaming under matched inputs and fixed tuned configurations. The complete protocol is provided in Supplementary Methods S1, and the paired results are reported in Supplementary Table S8.}

\subsection{Statistical analysis}

\TrackedChange{All Optuna comparisons use matched pairs within dataset, seed, and split units to control for substantial between-run heterogeneity. Within each matched unit, trials were ranked by the target metric, the top five were retained, and each condition was summarized by the median of those retained trials, yielding a top-$k$ summary with $k=5$ intended to estimate near-optimal attainable performance under a fixed number of tuning trials, where each trial is one candidate hyperparameter configuration evaluated by Optuna. Paired effects were then computed as simple condition differences, with negative values favoring the first condition for lower-is-better metrics. Dataset-level and global summaries are shown as forest plots with 95\% confidence intervals from two-sided paired one-sample $t$-tests. Supplementary Fig. S2 repeats the topology comparisons at $k\in\{1,3,5,10\}$, always using the median of the retained trials within each matched unit, to assess sensitivity to the number of near-optimal trials retained.}

We report Benjamini-Hochberg q-values alongside raw p-values for related dataset-level tests, and use q-values for figure markers and significance counts. Each topology contrast in Figs. 6-7 forms a separate family of 42 tests (14 datasets across $QE_B$, $QE_H$, and $QE_T$). Pooled overall tests are reported separately with raw p-values.

\section{Results}

\subsection{External SOM quality baselines}

\subsubsection{XPySOM calibration}

\TrackedChange{Under matched-configuration XPySOM versus FloatSOM calibration on hexagonal $QE$ (Fig. S1), no implementation-associated $QE$ difference was detected. Accordingly, we use hexagonal FloatSOM batch as the operational XPySOM reference in the benchmarks that follow, while not claiming formal statistical equivalence. The runtime differences are attributable to FloatSOM's JIT kernels, which incur a small startup cost. This overhead is progressively amortized as workload size increases, after which FloatSOM runs faster than XPySOM on larger datasets.}

\subsubsection{aweSOM quality baseline}

\TrackedChange{Relative to default aweSOM, untuned hexagonal FloatSOM reduced quantization error across all five evaluated datasets. Averaged equally across datasets, the improvements were 36.77\% for balanced $QE$, 28.52\% for holdout $QE$, and 45.11\% for train $QE$. The complete dataset-level results are reported in Supplementary Table S12. Given the large difference between aweSOM's serial-online training and the batch training used by FloatSOM and XPySOM, together with the large-workload timeout documented in the speed-benchmark protocol (Section 4.2), we used XPySOM as the external implementation comparator in subsequent analyses.}

\subsection{Comparison of Different Sampling Methods}

We first report a focused HDSSSOM pilot as an elimination comparison rather than as part of the broader sampling benchmark. This pilot used a smaller, more limited configuration than the later full versus random analysis and is summarized in Supplementary Table S2.

Under that pilot configuration, HDSSSOM was materially worse than full sampling. In the hexagonal view shown in Fig. 3, full outperformed HDSSSOM in all 50 paired comparisons (10 datasets $\times$ 5 seeds; 50 wins, 0 losses, 0 ties), with dataset-level median balanced-$QE$ improvements ranging from 2.5\% to 209.7\% and a global median improvement of 38.7\%. These HDSSSOM results should therefore be interpreted as a pilot screen under the stated configuration rather than as a comprehensive evaluation of all possible HDSSSOM schedules. Given these large differences, we focused on the full and random sampling methodologies for the remainder of the manuscript.

\begin{figure}[tb]
\centering
% SVG source: assets/figures/fig_3.svg
\includegraphics[width=\linewidth]{assets/figures/fig_3.pdf}
\caption{HDSSSOM screening pilot on $QE_B$ (hexagonal topology): full vs HDSSSOM, using the smaller pilot configuration summarized in Supplementary Table S2 (10 datasets, 5 seeds, with all other pilot settings held fixed). Panels report dataset-matched paired top-$k$ within-unit medians plus dataset-level paired-effect summaries (Section 4.4). Forest whiskers denote 95\% paired $t$-test confidence intervals around the mean paired effect.}
\label{fig:3}
\end{figure}

The full versus random analysis in Fig. 4 was generated from matched hexagonal Optuna runs in which the sampling selector was switched from full to random. Effects were computed within matched seed-specific units and then summarized across seeds. Fig. 4A-C summarize the matched full versus random paired effects for balanced, holdout, and train $QE$, respectively, showing that full sampling generally provides equal or better $QE$ than random sampling. Notably, this full versus random separation is much smaller than the full versus HDSSSOM pilot effect shown in Fig. 3; in most dataset-level comparisons, the full versus HDSSSOM improvement is at least twice as large - as evidenced by the difference in axis scales between Fig. 3 and Fig. 4.

We find the effectiveness of random sampling relative to full sampling is scale-dependent. Above $10{,}000$ samples, paired $QE$ differences are not meaningfully detected, whereas in smaller datasets the random arm shows higher variability and less stable outcomes (Fig. 4D-F). This is consistent with reduced per-iteration sample support under random subsampling (Fig. 4D-F). In this benchmark, the $>10{,}000$ regime is therefore a useful practical proxy for more stable random sampling behavior.

\begin{figure}[tb]
\centering
% SVG source: assets/figures/fig_4.svg
\includegraphics[width=\linewidth]{assets/figures/fig_4.pdf}
\caption{Full versus random sampling under the hexagonal topology. Panels A-C report paired full-versus-random effects for $QE_B$, $QE_H$, and $QE_T$. Panels D-F show the corresponding dataset-size regressions for the same three metrics, with numbered datasets keyed in Supplementary Table S1. Figure significance markers use Benjamini-Hochberg adjusted q-values. In the larger-dataset regime observed here (>10,000 samples), little paired $QE$ separation is detected; at smaller dataset scales, random is more variable and less stable.}
\label{fig:4}
\end{figure}

Nevertheless, full sampling remains the best sampling strategy for optimal $QE$ results. Accordingly, all remaining analyses reported below rely on full sampling unless specified.

\subsection{Topology Results}

Topology comparisons use $QE$ as the optimized endpoint and hexagonal batch SOM as the regular-topology baseline; Fig. 5 illustrates the three neighborhood structures. \texttt{initial\_radius} was optimized over the same interval for every topology. The selected full-sampling values were 1.03 for hexagonal, 1.46 for MST, and 1.41 for RNG; under random sampling they were 1.17, 1.82, and 1.77. Thus, the graph-topology gains are not comparisons against a broader default-radius hexagonal map, although radius was not isolated from the other tuned parameters.

\begin{figure}[tb]
\centering
% SVG source: assets/figures/fig_5.svg
\includegraphics[width=\linewidth]{assets/figures/fig_5.pdf}
\caption{Representative node and connection overlays for matched 100-node hexagonal, MST, and RNG SOMs. A--C: circles in native 2D. D--F: KDD Cup 99, trained on standardized 41-dimensional observations and displayed using one shared 2D PCA projection fitted to the observations and applied to all three sets of prototypes. Training used seed 42, full sampling, random initialization, 50 iterations, and the untuned XPySOM-like settings used at this point in the manuscript: initial learning rate 0.5, initial radius 5 with exponential decay, momentum disabled, and XPySOM-compatible normalization. Grey observation clouds show a deterministic maximum of 30,000 points; black lines are topology connections and red markers are SOM nodes. PCA can distort graph geometry in the original 41-dimensional space.}
\label{fig:5}
\end{figure}

In the KDD Cup 99 projection (Fig. 5D--F), the hexagonal topology contains several apparently nonlocal connections spanning separated regions of the projected prototype distribution, whereas the MST and RNG connections more closely follow its local geometry. This pattern illustrates how predetermined lattice neighbors can couple the updates of prototypes that are not locally adjacent in the displayed data structure, restricting how independently the nodes can redistribute.

\subsubsection{MST}

We report dataset-wise paired improvement summaries (hexagonal over MST) using the statistical procedure in Section 4.4 in Fig. 6.

\TrackedChange{Overall, MST outperforms matched hexagonal on Balanced QE (Fig. 6A), indicating a net advantage across train and holdout performance. This aggregate gain is driven more clearly by Train QE (Fig. 6C), while Holdout QE is more mixed across datasets (Fig. 6B) and shows no clear overall holdout advantage. The overall paired t-test p-values are Balanced QE (p=1.12e-05); Holdout QE (p=0.15); and Train QE (p=0.0064). Supplementary Table S5 lists the corresponding per-dataset and overall hexagonal-comparison effect estimates, 95\% confidence intervals, raw p-values, and Benjamini-Hochberg q-values for MST and RNG.}

\begin{figure}[tb]
\centering
% SVG source: assets/figures/fig_6.svg
\includegraphics[width=\linewidth]{assets/figures/fig_6.pdf}
\caption{\TrackedChange{Hexagonal versus MST topology on $QE$ metrics under full sampling only. Panels A-C report paired full sampling only $QE$ effects for $QE_B$, $QE_H$, and $QE_T$ across the available full sampling datasets. Forest whiskers denote 95\% paired $t$-test confidence intervals around the mean paired effect. Supplementary Table S5 reports the per-dataset effect estimates, 95\% confidence intervals, raw p-values, and dataset-level Benjamini-Hochberg adjusted q-values.}}
\label{fig:6}
\end{figure}

\subsubsection{RNG}

To evaluate RNG topology performance, we reuse the paired reporting logic on hexagonal versus RNG, again centered on $QE_B$ with $QE_H$ and $QE_T$ in Fig. 7.

\TrackedChange{RNG has lower QE than matched hexagonal on the reported QE endpoints (Fig. 7A-7C), with overall paired t-test p-values of Balanced QE (p=7.4e-10); Holdout QE (p=0.0232); and Train QE (p=4.69e-06). Supplementary Table S5 lists the corresponding per-dataset and overall hexagonal-comparison effect estimates, 95\% confidence intervals, raw p-values, and Benjamini-Hochberg q-values for MST and RNG.}

The main trend in Fig. 7 is that RNG improves on hexagonal most clearly in balanced QE and especially in train QE, with the separation most apparent in the real and larger datasets where the added flexibility of the graph neighborhood appears more useful than the fixed regular lattice.

\begin{figure}[tb]
\centering
% SVG source: assets/figures/fig_7.svg
\includegraphics[width=\linewidth]{assets/figures/fig_7.pdf}
\caption{\TrackedChange{Hexagonal versus RNG topology on $QE$ metrics under full sampling only. Panels A-C report paired full sampling only $QE$ effects for $QE_B$, $QE_H$, and $QE_T$ across the available full sampling datasets. Forest whiskers denote 95\% paired $t$-test confidence intervals around the mean paired effect. Supplementary Table S5 reports the per-dataset effect estimates, 95\% confidence intervals, raw p-values, and dataset-level Benjamini-Hochberg adjusted q-values.}}
\label{fig:7}
\end{figure}

\subsubsection{Direct MST--RNG comparison}

Fig. 8 directly compares MST and RNG under full sampling. In the Optuna matched top-$k$ runs, RNG achieved significantly lower balanced and train $QE$ than MST, while no significant difference was detected for holdout $QE$. Across the topology comparisons in Figs. 6--8, RNG therefore showed the highest quantization capacity among the evaluated topologies.

\begin{figure}[tb]
\centering
% SVG source: assets/figures/fig_8.svg
\includegraphics[width=\linewidth]{assets/figures/fig_8.pdf}
\caption{MST versus RNG topology on $QE$ metrics under full sampling only. Panels A-C report paired full sampling only $QE$ effects for $QE_B$, $QE_H$, and $QE_T$ across the available full sampling datasets. Forest whiskers denote 95\% paired $t$-test confidence intervals around the mean paired effect.}
\label{fig:8}
\end{figure}

\subsubsection{Initial-radius sensitivity}

Because initial radius was selected during tuning and differed across topology families, we isolated its effect in a fixed-configuration sweep. All non-radius hyperparameters were held at the tuned hexagonal full-sampling configuration, and each of seven radii was evaluated for hexagonal, MST, and RNG maps using the same 20 seeds on all 14 datasets. Fig. 9A reports $QE_B$ after normalizing every dataset--seed condition to its matched hexagonal value at the selected radius ($r=1.027$). We summarize these ratios by averaging log ratios across seeds within each dataset, weighting the 14 datasets equally, and exponentiating the cross-dataset mean. The resulting geometric-mean ratio is dimensionless: 1 denotes matched hexagonal performance at the selected radius, values below 1 indicate lower $QE_B$, and values above 1 indicate higher $QE_B$.

\TrackedChange{The exact normalized $QE_B$ radius-response values are reported in Supplementary Table S9. Separation vanished at $r\leq0.75$, with no significant MST--hexagonal or RNG--hexagonal difference at either $r=0.5$ or $r=0.75$ (Supplementary Table S10). The per-topology minimum $QE_B$ across the sweep occurred at $r=0.75$ for hexagonal and $r=1.5$ for MST and RNG. In paired tests between those minima, MST was 0.98\% lower than hexagonal (ratio 0.9902, 95\% CI [0.9747, 1.0060]; $p=0.203$, $q=0.304$), and RNG was 1.35\% lower than hexagonal (ratio 0.9865, 95\% CI [0.9689, 1.0043]; $p=0.125$, $q=0.304$); neither difference was significant. RNG and MST also did not differ significantly (ratio 0.9962, 95\% CI [0.9834, 1.0092]; $p=q=0.538$; Supplementary Table S11).}

Increasing radius eventually exchanged higher $QE_B$ for lower observed MTR in every topology, but this trade-off emerged earlier and much more sharply for hexagonal maps. RNG maintained the lowest observed MTR throughout the sweep, while the node-utilization curves did not indicate poorer use of the available map nodes. Because these raw MTR curves compare a fixed lattice with weight-derived graphs, their cross-family separation is descriptive. The significance symbols in Fig. 9 represent within-topology radius comparisons.

\begin{figure}[tb]
\centering
% SVG source: assets/figures/fig_9.svg
\includegraphics[width=\linewidth]{assets/figures/fig_9.pdf}
\caption{\TrackedChange{Initial-radius sensitivity under matched fixed configurations. A: dataset-balanced geometric-mean $QE_B$ ratio relative to the matched hexagonal configuration at its selected radius ($r=1.027$); lower values are better and the dashed horizontal line marks a ratio of 1. B: observed balanced Mean Tied Rank (MTR; lower is better). C: observed balanced node utilization (higher is better). All panels summarize 14 datasets with 20 matched seeds per topology and radius. The dotted vertical line and \texttt{Hex optimal} tick mark $r=1.027$. Significance markers use Benjamini--Hochberg adjusted tests comparing each non-anchor radius with $r=1.027$ within topology across the 54 radius--topology--metric tests. Confidence intervals are omitted from the plotted panels for readability; Supplementary Table S9 reports the exact normalized $QE_B$ values and intervals underlying panel A, and Supplementary Table S10 reports the post hoc matched cross-topology estimates and adjusted tests at every radius.}}
\label{fig:9}
\end{figure}

The fixed-configuration diagnostics provide the corresponding deployable tuned-default comparison (Table \ref{tab:matched_topology_diagnostics_summary}). Unlike the Optuna matched top-$k$ results, no balanced-$QE$ difference was detected between the tuned RNG and MST defaults ($p=0.126$). Within these two weight-derived graph families, RNG lowered mean balanced MTR from 7.87 to 4.91, a difference of 2.96 tied-rank positions ($p=1.75\times10^{-65}$), and increased balanced node utilization by 0.0040 ($p=6.07\times10^{-4}$).

\TrackedChange{The permutation-null rerun gave empirical null means between 49.999 and 50.012 across the six profile--topology summaries, consistent with the exact expectation of 50. Because the null values were effectively identical across topology families, graph--hexagonal MTR comparisons are reported descriptively, and MTR is used inferentially only to compare MST and RNG. Within these two weight-derived graph families, RNG showed significantly better local topological ordering than MST, as measured by null-relative MTR, in both the untuned (difference favouring RNG 0.0520, 95\% CI [0.0468, 0.0573], $p=1.33\times10^{-54}$) and tuned profiles (0.0614, 95\% CI [0.0561, 0.0668], $p=1.63\times10^{-65}$; Supplementary Table S13).}

\begin{table}[tb]
\caption{\TrackedChange{Matched topology diagnostics for balanced quantization error and Mean Tied Rank. Positive paired effects favor the second topology in each contrast after applying metric directionality; lower raw values are better for both $QE_B$ and $MTR_B$. Graph--hexagonal MTR effects are descriptive; MTR is used inferentially only for the RNG--MST comparison between the two weight-derived graph families. Asterisks mark prespecified pooled paired effects with raw $p<0.05$, except that graph--hexagonal MTR cells are unmarked; these tests are distinct from the dataset-level multiplicity-adjusted families. Fixed-adjacency permutation-null ratios and matched contrasts are reported in Supplementary Table S13. Full confidence intervals, applicable raw paired-test p-values, split-specific metrics, node utilization, and dead-node fraction are reported in Supplementary Table S7.}}
\label{tab:matched_topology_diagnostics_summary}
\begin{center}
\begin{tabular}{lrrrrrr}
\toprule
& \multicolumn{3}{c}{$QE_B$} & \multicolumn{3}{c}{$MTR_B$} \\
\cmidrule(lr){2-4}\cmidrule(lr){5-7}
Profile & MST vs hex & RNG vs hex & RNG vs MST & MST vs hex & RNG vs hex & RNG vs MST \\
\midrule
Untuned & +0.269\textsuperscript{*} & +0.249\textsuperscript{*} & -0.020\textsuperscript{*} & -0.20 & +2.44 & +2.64\textsuperscript{*} \\
Tuned & +0.064\textsuperscript{*} & +0.052\textsuperscript{*} & -0.012 & +22.71 & +25.67 & +2.96\textsuperscript{*} \\
\bottomrule
\end{tabular}
\end{center}
\end{table}

\subsection{Hyperparameter Tuning and Stability}

\subsubsection{Performance Gains from Hyperparameter Tuning}

Having shown that the graph-topology effects persist when radius is varied under a common tuned non-radius configuration, we next quantify the broader effect of complete topology-specific tuning relative to the untuned reference.

\TrackedChange{Fig. 10 compares the fixed tuned and untuned configurations using pooled topology-level $QE$ summaries; topology-stratified results are provided in Supplementary Fig. S3.}

\begin{figure}[tb]
\centering
% SVG source: assets/figures/fig_10.svg
\includegraphics[width=\linewidth]{assets/figures/fig_10.pdf}
\caption{Tuned configuration versus untuned reference $QE$ comparison across $QE_B$, $QE_H$, and $QE_T$, pooled across all topology runs under the matched pairing keys. Positive values indicate the tuned configuration outperforms the untuned reference; the global overall row pools all matched tuned configuration/untuned reference pairs across datasets. Forest whiskers denote 95\% paired $t$-test confidence intervals around the mean paired effect.}
\label{fig:10}
\end{figure}

The fixed tuned configurations had lower values for all three pooled $QE$ metrics than the untuned references, while observed balanced MTR increased by 24.34 tied-rank positions for hexagonal maps, 1.43 for MST, and 1.11 for RNG. These raw changes are reported descriptively: because the MTR scale and adjacency construction differ across topology families, they do not establish that tuning caused a larger loss of local ordering in the fixed lattice. The changes were not accompanied by lower node utilization.

The $QE$ improvement was observed in every topology family.

\subsubsection{Hyperparameter Stability Across Topology and Sampling}

To assess whether the derived hyperparameters are robust across runs and datasets, we compare the within-topology seed-to-seed tuned-parameter drift (Section 4.3.2). Overall, MST and RNG exhibit less variability than hexagonal maps under both full and random sampling (Fig. 11A-B). Across these panels, full sampling is generally the more stable setting. Mirroring the sample-size-dependent $QE$ performance, Fig. 11C shows that random-sampling hyperparameter stability also improves with increasing sample size.

\begin{figure}[tb]
\centering
% SVG source: assets/figures/fig_11.svg
\includegraphics[width=\linewidth]{assets/figures/fig_11.pdf}
\caption{Hyperparameter stability by sampling mode. A: selected parameter stability under full sampling for hexagonal, MST, and RNG topologies (lower stability score is better). B: selected parameter stability under random sampling for the same topologies. C: dataset-size stability regression under random sampling, using the selected parameter stability score against sample size (log10) across the included topology families.}
\label{fig:11}
\end{figure}

\subsection{Execution-Path Diagnostic}

\TrackedChange{The quality analyses above used the in-memory pathway. Repeating the tuned diagnostics through Ray streaming with matched data and configurations produced small differences, none of which reached significance before or after Benjamini-Hochberg correction (Supplementary Table S8). The remaining differences may reflect floating-point accumulation order. The following speed benchmarks use the Ray pathway.}

\section{Speed Scaling}

We next examine three aspects of FloatSOM runtime performance: the cost of full relative to random sampling, the scaling behavior obtained with additional GPUs, and the runtime implications of MST and RNG relative to a hexagonal topology.

\subsection{Random versus Full Sampling Runtime}

We report sample scaling runtime comparisons for full versus random sampling using the synthetic datasets as outlined in Section 4.2.1. We stratify by hexagonal, MST, and RNG, for 1, 2, and 4 GPUs (Fig. 12).

\begin{figure}[tb]
\centering
% SVG source: assets/figures/fig_12.svg
\includegraphics[width=\linewidth]{assets/figures/fig_12.pdf}
\caption{Sample scaling runtime comparison of full versus random sampling across $G\in\{1,2,4\}$ GPUs (A,B,C). Curves report mean wall-clock training time (s) under harmonized settings; error bars denote $\pm 1$ standard deviation across $n=3$ repeated runs per configuration. Color encodes topology (hexagonal, MST, RNG), and line style encodes sampling mode (full vs. random). Shaded x-axis regions indicate sample-size ranges that could not be run in that panel relative to the shared axis maximum due to timeouts. Lower values indicate faster execution.}
\label{fig:12}
\end{figure}

Across topologies, random sampling is faster than full sampling across all 1-, 2-, and 4-GPU comparisons, with similar proportional reductions at a given dataset size. With more GPUs, larger datasets can also be processed before timing out, although the runtime advantage of random over full narrows at the largest sample counts.

For the 1-GPU random sampling runs, the last successful 1-GPU run was at 100M samples, the same number of samples as run on full samples, despite random being much faster. We reason that the failed subsequent runs on larger samples were due to the transition from RAM operation to disk-backed operation. Even in random sampling mode, the implementation still stages the full dataset initially and loads full chunks before the subsampler selects the training samples. Therefore, the relevant cost is no longer only the reduced number of selected samples. We therefore treat the 1-GPU random failure at 500M samples as a disk-mode systems limitation rather than as evidence against the general random versus full runtime ordering.

\subsection{Multi-GPU topology scaling and Larger-Than-Memory context}

To examine parallel scaling, we benchmarked FloatSOM under full sampling across $G\in\{1,2,4,8\}$ GPUs using workload-scaling benchmarks that extend from standard in-memory settings to larger workloads that exceed available memory.

\begin{figure}[tb]
\centering
% SVG source: assets/figures/fig_13.svg
\includegraphics[width=\linewidth]{assets/figures/fig_13.pdf}
\caption{Multi-GPU scaling across $G\in\{1,2,4,8\}$ GPUs. Panels A-C show runtime (s) for dimension, sample, and grid size scaling workloads, respectively. Panels D-F show scaling efficiency for the same workloads, computed from the single-GPU baseline and the corresponding $G$-GPU runtime. Runtime error bars denote $\pm 1$ standard deviation across $n=3$ repeated runs per configuration; the 100\% efficiency reference line indicates ideal linear scaling.}
\label{fig:13}
\end{figure}

\subsubsection{GPU Scaling and Larger-Than-Memory Runtime Performance}

Fig. 13 shows that increasing GPU count improves performance in the sample scaling regime by increasing parallel throughput and delaying the transition to disk-backed execution. In the sample scaling benchmark, the 500,000,000 sample dataset requires disk backing under the 2-GPU configuration, whereas the 8-GPU configuration remains in RAM mode until the 1,000,000,000 sample dataset; when staging is still required, the disk-to-GPU path is distributed across more nodes. The 8-GPU RNG configuration processes 1,000,000,000 samples in 369.41 s (6.16 min). Note that this speed includes the time required to transfer data from shared storage to node-local shards, alongside the overhead associated with operating across multiple HPC nodes.

Grid-size scaling is the main exception: at the largest tested grid size (64), runtime decreases from 934.01 s (15.57 min) on 1 GPU to 880.83 s (14.68 min) on 8 GPUs, a 5.69\% reduction. Once map-size and topology-refresh costs dominate, additional GPUs therefore contribute little further speedup.

\subsubsection{GPU Scaling Efficiency}

We next consider GPU efficiency under strong scaling, relative to ideal linear scaling. Efficiency values with directly measured 1-GPU baselines summarize observed strong scaling. Values using extrapolated denominators are shown as descriptive estimates and are not used to establish a general efficiency improvement across workload sizes. Efficiencies above 100\% may reflect both parallel throughput and a changed memory/data-staging regime, rather than superlinear computation. The absolute runtime curves and directly measured-baseline comparisons therefore provide the primary scaling evidence.

\subsection{Topology Runtime Comparisons}

With that systems context in place, we next compare topology runtimes across hexagonal, MST, and RNG configurations on 8 GPUs (Fig. 14).

\begin{figure}[tb]
\centering
% SVG source: assets/figures/fig_14.svg
\includegraphics[width=\linewidth]{assets/figures/fig_14.pdf}
\caption{\TrackedChange{Topology runtime comparison at fixed $G=8$ GPUs under full sampling. Panels A-C report mean wall-clock runtime (s) for dimension, sample, and grid size scaling workloads, respectively, with topology traces for hexagonal, MST, and RNG. Error bars denote $\pm 1$ standard deviation across $n=3$ repeated runs per configuration. Detailed numerical summaries are available in the accompanying Zenodo archive.}}
\label{fig:14}
\end{figure}

In Fig. 14A-B, the topologies scale similarly as input complexity and data volume increase: even at the largest tested axis values, the maximum pairwise runtime spread remains modest at dimension scaling (4.70\% at 5,000 dimensions) and sample scaling (3.26\% at 1,000,000,000 samples).

However, when the grid itself is enlarged in Fig. 14C, topology-dependent runtime differences become readily evident. At the largest tested grid size (grid size 64), the 8-GPU mean runtimes are 32.54 s (0.54 min) for hexagonal, 266.45 s (4.44 min) for MST, and 880.83 s (14.68 min) for RNG, corresponding to 8-GPU MST and RNG runtimes that are 8.19x and 27.07x the hexagonal runtime, respectively.

To further examine the RNG runtime spike in grid-size scaling, we compared the topology-construction operations. For each of the $\binom{P}{2}$ candidate node pairs, RNG checks every possible third node in its relative-neighborhood blocker test, giving up to $\binom{P}{2}(P-2)=O(P^3)$ pair--blocker comparisons per topology refresh. At $P=4096$, this is approximately 34.3 billion unordered pair--blocker checks. MST instead constructs and sorts approximately $\binom{4096}{2}=8.39$ million candidate edges using Kruskal's algorithm, an $O(P^2\log P)$ edge-construction step after pairwise distances are calculated. Both topologies then perform the shared graph-distance and influence calculations. The additional $O(P^3)$ RNG blocker-test work therefore provides a quantitative algorithmic explanation for the sharp RNG cost increase at large grid sizes.

\section{Final FloatSOM RNG Comparison with XPySOM}

\TrackedChange{Fig. 15 compares untuned hexagonal XPySOM with tuned FloatSOM RNG and therefore combines implementation, tuning, and topology effects \citep{manciniXPySomHighPerformanceSelfOrganizing2020}. Sections 5.1, 5.3, and 5.4 separate these components; the corresponding detailed tuned-hexagonal and tuned-MST comparisons are available in the accompanying Zenodo archive.}

\begin{figure}[tb]
\centering
% SVG source: assets/figures/fig_15.svg
\includegraphics[width=\linewidth]{assets/figures/fig_15.pdf}
\caption{Integrated deployment comparison of untuned hexagonal XPySOM versus tuned FloatSOM RNG, combining implementation, tuning, and topology effects. Panels A-C compare $QE_B$, $QE_H$, and $QE_T$ under full sampling; panel D provides the scaling/runtime context.}
\label{fig:15}
\end{figure}

At the overall level, Fig. 15 shows median percentage improvements of $QE_B$ (14.5\%); $QE_H$ (9.1\%); and $QE_T$ (22.5\%) for tuned FloatSOM RNG relative to untuned hexagonal XPySOM, capturing the combined deployment effect of implementation, topology choice, and tuning on $QE$.

\TrackedChange{The matched untuned hexagonal implementation comparison showed no detectable global balanced-$QE$ gain (median improvement 0.000016\%, $p=0.0923$; Supplementary Table S4); on the common untuned-XPySOM reference scale, tuning accounted for 13.04 percentage points of the full tuned-RNG stack's 14.49\% gain, and topology accounted for the remaining 1.44 points (0.88 for holdout $QE$ and 2.34 for train $QE$; Supplementary Table S6 and the accompanying Zenodo archive), consistent with the direct 1.297\% tuned RNG-over-hexagonal balanced-$QE$ advantage in Supplementary Table S5.}

For the untuned hexagonal XPySOM reference in Fig. 15, workloads beyond the $10^8$-sample case were not processed because they exceeded available VRAM and XPySOM requires the full dataset to be loaded into memory. Taken together, Fig. 15 shows the point at which the workload size is large enough to warrant the extra startup overhead incurred by tuned FloatSOM RNG: tuned FloatSOM RNG delivers better $QE$ than the untuned hexagonal XPySOM baseline, while also running faster and scaling to larger workloads.

\section{Discussion}

\TrackedChange{In this paper, we introduce FloatSOM, a distributed, out-of-memory GPU framework for large-scale SOM training that supports training-time graph topologies. FloatSOM scales SOM training across multiple GPUs and HPC nodes and supports datasets larger than available device memory through disk-backed streaming.}

The trade-off between full and random sampling is scale dependent. Random sampling is less stable on small datasets, while above $10{,}000$ samples we detect little paired $QE$ difference. Full sampling is therefore preferable when stability is critical; random sampling offers higher throughput on larger datasets, although disk-backed operation limits its I/O advantage.

Across the Optuna matched top-$k$ runs, RNG showed significantly better balanced and train $QE$ than MST, while no significant holdout-$QE$ difference was detected. RNG also had better overall hyperparameter stability. In the tuned defaults, no balanced-$QE$ difference was detected, although RNG had higher node utilization than MST.

A geometric interpretation consistent with these results is that a fixed lattice imposes uniform, predetermined coupling: when a prototype moves, it can unnecessarily influence lattice neighbors that are unrelated in the learned data space, potentially displacing useful prototypes and increasing dead nodes. The apparently nonlocal hexagonal connections in the KDD Cup 99 projection (Fig. 5D) provide a visual illustration of this coupling and the associated restriction on independent prototype redistribution, subject to the distortions inherent in the two-dimensional PCA display. MST minimizes such coupling, allowing prototypes to redistribute more freely along irregular or elongated data structures. This freedom is also its limitation, because a tree cannot retain multiple locally appropriate connections where a concentrated region is better represented by a mesh. RNG occupies an intermediate position. Its nodes can influence several neighbors, but those connections arise from the evolving prototype geometry rather than being imposed before training; RNG can remain sparse where appropriate and form multiple connections in dense regions. Within the two weight-derived graph families, RNG's higher node utilization is consistent with this interpretation \citep{kohonenEssentialsSelforganizingMap2013,kangasVariantsSelforganizingMaps1990,toussaintRelativeNeighbourhoodGraph1980}.

\TrackedChange{Among the graph topologies evaluated, RNG performed best on $QE$, stability, and node utilization, though its improvement over a tuned hexagonal lattice was modest. Where a graph topology is used at larger node counts, MST is preferable due to its lower compute cost.}

The quality experiments use the in-memory pathway, whereas the scaling experiments use Ray. Within the distributed benchmark, scaling depends on workload size: overhead dominates small problems, while larger workloads benefit from parallel throughput and may remain in RAM when the 1-GPU case requires disk backing. Efficiencies above 100\% reflect this change in memory regime rather than super-linear computation.

Regarding existing distributed SOM implementations, Somoclu and GigaSOM are important parallel-systems references. However, Somoclu was not benchmarked against because both Somoclu and SomocluGPU were already compared directly with XPySOM in the XPySOM study and found to be more than 10 times slower \citep{manciniXPySomHighPerformanceSelfOrganizing2020}. GigaSOM is a CPU-based algorithm implemented in Julia with perhaps the strongest demonstrated speed and size scale-ups. However, considering GigaSOM's hardware requirements, programming language difference, and dataset availability, a full systems benchmark was outside the scope of this paper. For context though, in GigaSOM's large dataset example completed a workflow that trained a 1024 node SOM on 1,167,129,317 events with 18 features in $<25$ minutes on an 11-node, 256-core CPU cluster \citep{kratochvilGigaSOMjlHighperformanceClustering2020}. Comparatively, FloatSOM trained a same-sized 1024-node SOM on 1,000,000,000 samples with 50 features, a dataset \textasciitilde{}x2.4 larger than GigaSOM's, in 6.16 minutes using 8 GPUs across two HPC nodes.

\TrackedChange{At grid size 64, hexagonal, MST, and RNG required 32.54, 266.45, and 880.83 s on 8 GPUs, respectively. RNG's graph-construction and all-pairs path calculations therefore scale much more sharply with grid size. More generally, full sampling favors stability, random sampling favors throughput on large datasets, and additional GPUs are most useful when they keep the workload in RAM.}

FloatSOM brings these topology and systems choices together in a single large-scale SOM framework.

% Acknowledgments and funding information are omitted from the anonymous
% review source and must be restored only after acceptance and deanonymization.

\bibliography{main}
\bibliographystyle{tmlr}
\clearpage

\appendix
\section{Supplementary Methods and Tables (End Matter)}

\subsection{Supplementary Methods S1. Execution-path concordance protocol}

The execution-path analysis used all 14 datasets, the 10 shared seeds 42--51, full sampling, and the hexagonal, MST, and RNG topologies on one NVIDIA V100 GPU. For each dataset--seed--topology unit, local CuPy and Ray streaming received the same standardized data, deterministic train--holdout split, topology-specific tuned parameters from Section 4.3.1, and evaluation metrics; execution profile was the only deliberately changed factor. Differences were defined as Ray streaming minus local CuPy.

\TrackedChange{For each topology, Supplementary Table S8 reports paired differences for balanced $QE$, MTR, node utilization, and dead-node fraction across 140 matched dataset--seed units. Mean differences and 95\% confidence intervals were calculated from two-sided one-sample paired $t$-tests against zero, with Benjamini--Hochberg correction applied jointly across the 12 topology--metric tests.}

\subsection{Supplementary Methods S2. MTR fixed-adjacency permutation-null protocol}

For every trained map and evaluation split, we held the final graph and the observed first--second BMU identity pairs fixed. Prototype identities were then uniformly permuted among the $P=100$ graph nodes 1,000 times, and MTR was recomputed after each permutation. MST and RNG graphs were not rebuilt. We calculated the empirical null mean and the observed-to-null ratio

\begin{equation*}
R_{\mathrm{MTR}}=
\frac{\mathrm{MTR}_{\mathrm{observed}}}
{\overline{\mathrm{MTR}}_{\mathrm{null}}},
\end{equation*}

where lower values indicate stronger ordering relative to random prototype placement within the fitted graph. Split-specific statistics were balanced by averaging the train and holdout values, as for the other diagnostic metrics. Matched topology contrasts were calculated across the 280 shared dataset--seed maps using two-sided one-sample paired $t$-tests.

For any connected $P$-node graph, the exact permutation-null expectation is $P/2$. Conditional on the first BMU, the permuted second BMU is uniform over the other $P-1$ nodes. When counted with their shell multiplicities, the tied ranks partition ordinal ranks 1 through $P-1$, whose mean is $P/2$. The theoretical null mean is therefore exactly 50 for the 100-node maps, irrespective of graph degree or shell sizes.

\subsection{Supplementary Tables}
\setcounter{table}{0}
\renewcommand{\tablename}{Supplementary Table}
\renewcommand{\thetable}{S\arabic{table}}
\renewcommand{\theHtable}{supplementarytable.\arabic{table}}

\input{tracked/supplementary_tables_v2_restructuring.tex}
\section{Supplementary Figures (End Matter)}
\setcounter{figure}{0}
\renewcommand{\figurename}{Supplementary Figure}
\renewcommand{\thefigure}{S\arabic{figure}}
\renewcommand{\theHfigure}{supplementaryfigure.\arabic{figure}}

\begin{figure}[t]
\centering
% SVG source: assets/figures/supp_fig_s1.svg
\includegraphics[width=\linewidth]{assets/figures/supp_fig_s1.pdf}
\caption{\TrackedChange{Hexagonal FloatSOM--XPySOM implementation calibration. Panels A-C report paired $QE$ effects for $QE_B$, $QE_H$, and $QE_T$. Panel D reports dataset-level median runtime deltas against dataset size, where each numbered dot is the median matched-seed value of \texttt{FloatSOM time - XPySOM time}; negative values favor FloatSOM and positive values favor XPySOM. The point numbers map to Supplementary Table S4. Forest whiskers denote exact distribution-free 95\% intervals for the paired Wilcoxon location estimate.}}
\label{fig:s1}
\end{figure}

\begin{figure}[t]
\centering
% SVG source: assets/figures/supp_fig_s2.svg
\includegraphics[width=\linewidth]{assets/figures/supp_fig_s2.pdf}
\caption{\TrackedChange{Topology sensitivity analyses under full sampling. A: hexagonal versus MST; B: hexagonal versus RNG; C: MST versus RNG. Within each row, columns report matched top-$k$ paired sensitivity analyses for $QE_B$, $QE_H$, and $QE_T$ at $k\in\{1,3,5,10\}$; directional labels indicate which topology is favored.}}
\label{fig:s2}
\end{figure}

\begin{figure}[t]
\centering
% SVG source: assets/figures/supp_fig_s3.svg
\includegraphics[width=\linewidth]{assets/figures/supp_fig_s3.pdf}
\caption{\TrackedChange{Tuned versus untuned fixed-configuration comparisons by topology. A: hexagonal; B: MST; C: RNG. Within each row, columns report $QE_B$, $QE_H$, and $QE_T$ for 14 datasets and 20 seeds under matched dataset--seed--topology--split keys. Positive values indicate that the tuned configuration outperforms the untuned reference; profile definitions are given in Section 4.3.1.}}
\label{fig:s3}
\end{figure}

\begin{figure}[t]
\centering
% SVG source: assets/figures/supp_fig_s4.svg
\includegraphics[width=\linewidth]{assets/figures/supp_fig_s4.pdf}
\caption{\TrackedChange{Full GPU-count scaling context for MST and hexagonal under matched full sampling settings. A-C: MST dimension, sample, and grid-size scaling; D-F: hexagonal dimension, sample, and grid-size scaling. Curves correspond to $G\in\{1,2,4,8\}$ GPUs and report mean wall-clock runtime (s) with $\pm 1$ standard-deviation error bars across $n=3$ repeated runs per configuration.}}
\label{fig:s4}
\end{figure}

\end{document}
